\documentclass[preview]{standalone}

\usepackage{amsmath}
\usepackage{amssymb}
\usepackage{parskip}
\usepackage{fullpage}
\usepackage{hyperref}
\usepackage{stellar}
\usepackage{definitions}
\usepackage{bettelini}

\begin{document}

\id{diffeq-laplace-transform}
\genpage

\section{Laplace Transform}

\begin{snippetdefinition}{laplace-transform-definition}{Laplace Transform}
    Let \(f\colon[0,+\infty)\fromto\realnumbers\) be piecewise continuous
    on every compact interval. Its \emph{Laplace transform}, at each real \(s\)
    for which the following improper integral converges absolutely, is
    \[
        \mathcal{L}\{f(t)\}= \integral[0][\infty][e^{-st}f(t)][t]=F(s)
    \]
\end{snippetdefinition}

\subsection{Properties of the Laplace Transform}

\begin{snippetcorollary}{laplace-transform-constant-property}{Laplace Transform constant property}
    Let \(f,g\) have absolutely convergent Laplace transforms at \(s\).
    Then
    \[
        \laplacetf\{af(t)+bg(t)\} = a\laplacetf\{f(t)\} + b\laplacetf\{g(t)\}
    \]
    for any constants \(a\) and \(b\).
\end{snippetcorollary}

\subsection{Inverse Laplace Transform}

\begin{snippetdefinition}{inverse-laplace-transform-definition}{Inverse Laplace Transform}
    On continuous \function[functions] on \([0,+\infty)\) of exponential order,
    the \emph{inverse Laplace transform} recovers \(f\) from its transform \(F\):
    \[
        {\mathcal{L}}^{-1} \{F(s)\}= f(t)
    \]
\end{snippetdefinition}

\subsection{Properties of the Inverse Laplace Transform}

\begin{snippetcorollary}{inverse-laplace-transform-contant-property}{Inverse Laplace Transform constant property}
    Let \(F,G\) be the transforms of continuous \function[functions] \(f,g\)
    of exponential order on \([0,+\infty)\). Then
    \[
        \invlaplacetf \{aF(s)+bG(s)\} =
        a\invlaplacetf\{F(s)\} +
        b\invlaplacetf\{G(s)\}
    \]
    for any constants \(a\) and \(b\).
\end{snippetcorollary}

\subsection{Heaviside function}

%%%%%%%%%%%%%%%%%%%%%%%%% FROM HERE

\begin{snippettheorem}{heaviside-function-expl}{Translation in time}
    Let \(c\geq0\) and let \(f\) be piecewise continuous on compact subintervals
    of \([0,+\infty)\), with absolutely convergent Laplace transform \(F(s)\).
    Define \(g(t)=0\) for \(t<c\), and \(g(t)=f(t-c)\) for \(t\geq c\).
    Equivalently \(g(t)=H(t-c)f(t-c)\), with the zero branch understood
    without evaluating \(f\) at negative arguments. Then
    \[
        \laplacetf\{g\}(s)=e^{-cs}F(s).
    \]
\end{snippettheorem}

\begin{snippetproof}{heaviside-function-expl-proof}{heaviside-function-expl}{Translation in time}

    \begin{align*}
        \laplacetf\{H(t-c)f(t-c)\} &=
        \integral[0][\infty][e^{-st}H(t-c)f(t-c)][t] \\
        &= \integral[c][\infty][e^{-st}f(t-c)][t]
    \end{align*}
    Substitute \(u=t-c\):
    \begin{align*}
        \integral[0][\infty][e^{-s(u+c)}f(u)][u]
        &= \integral[0][\infty][e^{-su}e^{-cs}f(u)][u] \\
        &= e^{-cs} \integral[0][\infty][e^{-su}f(u)][u]
    \end{align*}

    Concluding that

    \[
        \laplacetf\{H(t-c)f(t-c)\} =
        e^{-cs}F(s)
    \]
\end{snippetproof}

\begin{snippetnote}{diffeq-laplace-inversion-convention}{Values at jumps}
    For piecewise continuous \function[functions], a Laplace transform does not
    distinguish changes at isolated points. Inverse transforms in this class
    therefore specify the \function at its continuity points, with a separate
    convention at jumps. In this sense the translation formula is also written
    \(\invlaplacetf\{e^{-cs}F(s)\}=H(t-c)f(t-c)\).
\end{snippetnote}

\subsection{Laplace Transform of derivatives}

\begin{snippettheorem}{laplace-transform-of-derivatives}{Laplace Transform of derivatives}
    Let \(n\geq1\), \(f\in C^{n-1}([0,+\infty))\), and suppose
    \(f^{(n-1)}\) is piecewise continuously differentiable on compact intervals.
    Assume there exist \(M,a,T\) such that
    \(|f^{(j)}(t)|\leq Me^{at}\) for \(0\leq j\leq n\), \(t\geq T\),
    wherever the derivatives are defined. Then for real \(s>a\),
    \[
        \laplacetf\{f^{(n)}\} =
        s^n \laplacetf\{f\} -
        \sum_{k=1}^n s^{n-k}f^{(k-1)}(0)
    \]
\end{snippettheorem}

\subsection{Solving Initial value problems}

\begin{snippet}{solving-iv-problems}
Given an initial value problem differential equation, we can apply the Laplace
transform on both sides of the equation. After applying the initial conditions
(\(y(0), y'(0), \cdots\)), the differential equation is transformed into an algebraic equation.

Note that if the initial conditions are not expressed as  \(y(0), y'(0), \cdots\),
we need to first make a change of variable.

We can then isolate \(\laplacetf\{y\}\) in the equation and get
\[
    \laplacetf\{y\} = M
\]
Then apply the inverse Laplace transform to solve the equation
\[
    y=\invlaplacetf\{M\}
\]
\end{snippet}

% from https://tutorial.math.lamar.edu/Classes/DE/IVPWithNonConstantCoefficient.aspx


\end{document}
